\section{DISCUSSION, RESEARCH ROADMAP, AND LIMITATIONS}
\label{sec:discussion_roadmap}

Across the evidence base, the central issue is whether physical sharing
produces joint outcomes that remain traceable to one operating context. This
section interprets that relationship, develops five research priorities, and
states the review boundaries.

\subsection{Interpretation Across the Evidence Base}

Physical integration identifies where communication and sensing share
infrastructure, signals, or information. Evidential integration shows how
clearly both outcomes remain tied to the same operating conditions. Together,
they connect architectural and transmission perspectives while preserving the
physical regime of each platform
\cite{Wen2024OISACArchitectures,Su2025OISACTheoryPractice,
Liang2024LiSACReview,Liu2026OpticalFiberISACReview}.

Deeper sharing can connect more design variables and improve coordination
across the two functions. It can also couple power, distortion,
synchronization, dynamic range, and failure behavior. Its scientific value
therefore depends on locating these effects in the signal chain and tracing
their influence on both outcomes.

Tradeoff patterns transfer across studies when tasks, metric definitions,
measurement planes, baselines, constraints, and validation settings remain
aligned. This alignment allows model bounds and empirical operating points to
retain their distinct roles and support comparison within a stated context
\cite{Khorasgani2026GeneralOpticalISACSurvey,Rode2025PolarizationISAC}.

Fielded cable and network studies illustrate the value of operational
continuity
\cite{Cao2026SubmarineFiberISAC,Ip2026DeployedTelecomCableISAC}. Within the
survey, system evidence becomes reusable when configuration, calibration,
processing, and reported artifacts remain connected across interfaces. This
continuity links physical sharing to application and network claims and
motivates a central research task. Joint operating points should remain
interpretable, testable, and reusable across groups.

\subsection{Research Priorities for Reusable Evidence}

Table~\ref{tab:research_roadmap} organizes five priorities and pairs each
observed finding with an action and a measurable sign of progress.

\begin{table}[!t]
\caption{Research priorities derived from the reviewed evidence}
\label{tab:research_roadmap}
\centering
\footnotesize
\setlength{\tabcolsep}{2.4pt}
\renewcommand{\arraystretch}{1.02}
\begin{tabularx}{\columnwidth}{@{}
>{\raggedright\arraybackslash}p{0.20\columnwidth}
>{\raggedright\arraybackslash}X
>{\raggedright\arraybackslash}p{0.30\columnwidth}@{}}
\toprule
Priority & Finding and proposed action & Evidence of progress \\
\midrule
\textbf{1. Measurement contract} &
\emph{Finding.} Of 4,779 metric records, 118 supported bounded relations under
aligned conditions. \emph{Action.} Record the task, plane, unit, baseline,
operating condition, validation setting, and source locator. &
An independent reader reconstructs the quantity and condition from a structured
record. \\
\addlinespace[1.5pt]
\textbf{2. Benchmark for each modality} &
\emph{Finding.} Comparison across studies depends on matched physical context.
\emph{Action.} Repeat one task with reference and independent designs under
shared calibration and ground truth. &
Independent groups reproduce the relative finding under the shared protocol. \\
\addlinespace[1.5pt]
\textbf{3. Joint operation under disturbance} &
\emph{Finding.} Of 402 substantive tradeoff relationships, 371 depend on
reported conditions. Six of 12 field or deployment tier studies report
outcomes in both domains. Traceability is strongest when both outcomes share
the same timing and condition set.
\emph{Action.} Evaluate both functions in the same run and condition set under
load, a disturbance relevant to the modality, and matched baselines for the
individual functions. &
Repeated trials recover the joint region and its limiting mechanism. \\
\addlinespace[1.5pt]
\textbf{4. Reconstructable experiment package} &
\emph{Finding.} Open data appear in 13 of 206 studies, and code or model access
reaches a partial release in one. \emph{Action.} Preserve a versioned package
and rerun the main comparison under one added condition. &
An independent rerun recovers the comparison and added condition from the
package. \\
\addlinespace[1.5pt]
\textbf{5. Intelligence, scale, and security} &
\emph{Finding.} Technology and application coverage connects the corpus to 6G
roles. Conformance, interoperability, and readiness become measurable through
explicit network tests. \emph{Action.} Apply a declared shift, attack, load, or
component failure and report overhead, resource use, decisions, and fallback. &
The service retains an explainable boundary with resource use, decisions, and
fallback recorded. \\
\bottomrule
\end{tabularx}
\vspace{1pt}
\parbox{\columnwidth}{\raggedright\footnotesize\emph{Note.} The findings
summarize the retained corpus presented in Sections~\ref{sec:optical_modality_families}
through \ref{sec:technologies_applications_6g}. The proposed actions retain the
physical conditions of each modality and define observable signs of progress.\par}
\end{table}

The order matters because each stage uses evidence produced by the previous
one. Clear measurement context supports fair benchmarks. Benchmarks support
realistic joint tests. Reconstructable artifacts make those tests repeatable,
and repeatable tests provide a basis for adaptive and secure evaluation at
network scale.

Stress tests should follow the native signal path of each modality. Fiber tests
can vary topology and traffic, free space optical tests can vary geometry and
atmosphere, and visible light tests can vary illumination and orientation.
Photonic terahertz tests can vary optical drift together with wireless
propagation. Matched baselines for the individual functions keep the limiting
mechanism visible across these settings.

For standards and 6G coordination, the same sequence can support a core
reporting profile that preserves modality physics, ground truth, and reusable
benchmark definitions. Shared metadata can connect results across platforms
while retaining the conditions needed for interpretation.

\subsection{Scope and Limitations of the Review}

The review represents English language technical reports published from 1
January 2020 through 22 June 2026 and indexed by the six searched sources. Full
texts were available and assessed for 272 of the 330 reports sought, which
corresponds to 82.4\% retrieval coverage. Two Taylor \& Francis exports with few
records are documented through their complete platform exports. These
boundaries define the literature represented by the final corpus of 206 studies.

Searches and screening preceded retrospective registration, and the subsequent
search refresh and report reconciliation produced the final corpus described
in Section~\ref{sec:review_methods}. One investigator supervised routine
decisions, while independent checks focused on specified evidence records and
report reconciliation.

TQAF provides a technical appraisal developed for this review. The diversity
of tasks, quantities, validation settings, outcomes, and source protocols
supports structured narrative synthesis, which preserves native conditions and
bounded comparisons. This analytical scale centers interpretation on reported
mechanisms and conditions. Harmonized outcomes and source protocols form the
appropriate basis for pooled effects, sensitivity estimates, and causal
moderator analysis.

The multilabel taxonomy incorporates reviewer judgment, especially for hybrid
systems. Row counts represent coded observations at extraction granularity,
while effect analysis uses a different unit. For 92 metric rows, the available
fields supported interpretation within the source setting and a conservative
comparison status. Artifact access combined source reporting with targeted link
checks, so the appraisal reflects documented reconstruction information.

The conclusions concern recurring mechanisms, measurement constraints, and
validation patterns within the retained corpus. Corpus distributions describe
the reviewed literature. Worldwide prevalence, market maturity, and platform
rank require evidence designed for those questions. Within its stated scope,
the review identifies the conditions that allow optical integration to support
transferable joint results.
